\documentclass{amsart}
% For dealing with references we use the comment environment
\usepackage{verbatim}
\newenvironment{reference}{\comment}{\endcomment}
%\newenvironment{reference}{}{}
\newenvironment{slogan}{\comment}{\endcomment}
\newenvironment{history}{\comment}{\endcomment}

% For commutative diagrams we use Xy-pic
\usepackage[all]{xy}

% We use 2cell for 2-commutative diagrams.
\xyoption{2cell}
\UseAllTwocells

% We use multicol for the list of chapters between chapters
\usepackage{multicol}

% This is generally recommended for better output
\usepackage{lmodern}
\usepackage[T1]{fontenc}

% For cross-file-references

% Package for hypertext links:
\usepackage{hyperref}

% For any local file, say "hello.tex" you want to link to please

% Theorem environments.
%
\theoremstyle{plain}
\newtheorem{theorem}[subsection]{Theorem}
\newtheorem{proposition}[subsection]{Proposition}
\newtheorem{lemma}[subsection]{Lemma}

\theoremstyle{definition}
\newtheorem{definition}[subsection]{Definition}
\newtheorem{example}[subsection]{Example}
\newtheorem{exercise}[subsection]{Exercise}
\newtheorem{situation}[subsection]{Situation}

\theoremstyle{remark}
\newtheorem{remark}[subsection]{Remark}
\newtheorem{remarks}[subsection]{Remarks}

\numberwithin{equation}{subsection}

% Macros
%
\def\lim{\mathop{\mathrm{lim}}\nolimits}
\def\colim{\mathop{\mathrm{colim}}\nolimits}
\def\Spec{\mathop{\mathrm{Spec}}}
\def\Hom{\mathop{\mathrm{Hom}}\nolimits}
\def\Ext{\mathop{\mathrm{Ext}}\nolimits}
\def\SheafHom{\mathop{\mathcal{H}\!\mathit{om}}\nolimits}
\def\SheafExt{\mathop{\mathcal{E}\!\mathit{xt}}\nolimits}
\def\Sch{\mathit{Sch}}
\def\Mor{\mathop{\mathrm{Mor}}\nolimits}
\def\Ob{\mathop{\mathrm{Ob}}\nolimits}
\def\Sh{\mathop{\mathit{Sh}}\nolimits}
\def\NL{\mathop{N\!L}\nolimits}
\def\CH{\mathop{\mathrm{CH}}\nolimits}
\def\proetale{{pro\text{-}\acute{e}tale}}
\def\etale{{\acute{e}tale}}
\def\QCoh{\mathit{QCoh}}
\def\Ker{\mathop{\mathrm{Ker}}}
\def\Im{\mathop{\mathrm{Im}}}
\def\Coker{\mathop{\mathrm{Coker}}}
\def\Coim{\mathop{\mathrm{Coim}}}

% Boxtimes
%
\DeclareMathSymbol{\boxtimes}{\mathbin}{AMSa}{"02}

%
% Macros for moduli stacks/spaces
%
\def\QCohstack{\mathcal{QC}\!\mathit{oh}}
\def\Cohstack{\mathcal{C}\!\mathit{oh}}
\def\Spacesstack{\mathcal{S}\!\mathit{paces}}
\def\Quotfunctor{\mathrm{Quot}}
\def\Hilbfunctor{\mathrm{Hilb}}
\def\Curvesstack{\mathcal{C}\!\mathit{urves}}
\def\Polarizedstack{\mathcal{P}\!\mathit{olarized}}
\def\Complexesstack{\mathcal{C}\!\mathit{omplexes}}
% \Pic is the operator that assigns to X its picard group, usage \Pic(X)
% \Picardstack_{X/B} denotes the Picard stack of X over B
% \Picardfunctor_{X/B} denotes the Picard functor of X over B
\def\Pic{\mathop{\mathrm{Pic}}\nolimits}
\def\Picardstack{\mathcal{P}\!\mathit{ic}}
\def\Picardfunctor{\mathrm{Pic}}
\def\Deformationcategory{\mathcal{D}\!\mathit{ef}}

\setlength{\emergencystretch}{3em}
\begin{document}
\title[Stacks Project, Tag 0F3R]{566 --- The maximal weighted-fibre locus of a quasi-finite finitely presented affine morphism with positive weighting is affine}
\maketitle
\noindent Copyright (C) 2005--2025 Johan de Jong. Stacks Project authors. Source: \href{https://stacks.math.columbia.edu/tag/0F3R}{Tag 0F3R}.

\noindent Editorial difficulty note (not author text). Difficulty H2: The proof combines Noetherian approximation of the weighting with finite surjective affineness descent. After replacing the finite compactification by a union of sections, it computes how generic weights specialize and identifies the maximal locus as an intersection of affine section images..

\noindent Editorial provenance note (not author text). History: excerpt from revision \texttt{a0444\allowbreak 6e57e\allowbreak c1fbc\allowbreak 252a8\allowbreak 71afc\allowbreak ec775\allowbreak 2fb28\allowbreak 07b14}, more-morphisms.tex, lines 22965--23077. Reference presentation was adapted; original-author context and core support blocks were retained or added. The mathematical wording is unchanged. Licensed under GNU FDL 1.2 or later; no invariant sections or cover texts. See ../licenses/COPYING. Context blocks reproduce author definitions and supporting results; they are not additional counted problems.

\section{Weightings}
\label{section-weightings}

\noindent
The material in this section is taken from \href{https://stacks.math.columbia.edu/bibliography}{Bibliography: SGA4 (Exposee XVII, 6.2.4)}.

\medskip\noindent
Let $\pi : U \to V$ be a locally quasi-finite morphism of schemes with
finite fibres. Given a function $w : U \to \mathbf{Z}$ we define a function
$$
\textstyle{\int}_\pi w : V \longrightarrow \mathbf{Z},\quad
v \longmapsto
\sum\nolimits_{u \in U,\ \pi(u) = v} w(u) [\kappa(u) : \kappa(v)]_s
$$
Note that the field extensions are finite
(Morphisms, Lemma \href{https://stacks.math.columbia.edu/tag/01TG}{lemma residue field quasi finite (Tag 01TG)}),
$[\kappa' : \kappa]_s$ is the separable degree
(Fields, Definition \ref{fields-definition-insep-degree}),
and the sum is finite as the fibres of $\pi$ are assumed finite.
Another way to compute the value of $\int_\pi w$ at a point $v \in V$
is as follows.
Choose an algebraically closed field $k$ and a morphism
$\overline{v} : \Spec(k) \to V$ whose image is $v$. Then we have
$$
(\textstyle{\int}_\pi w)(v) =
\sum\nolimits_{\overline{u} \in U_{\overline{v}}} w(\overline{u})
$$
where of course $w(\overline{u})$ denotes the value of $w$ at the
image $u$ of the point $\overline{u}$ under the morphism
$U_{\overline{v}} \to U$.
Note that we may view $\overline{u} \in U_{\overline{v}}$ as morphisms
$\overline{u} : \Spec(k) \to U$ such that
$\pi \circ \overline{u} = \overline{v}$. Namely, since
$U \to V$ is locally quasi-finite with finite fibres,
the scheme $U_{\overline{v}}$ is the spectrum
of a finite dimension algebra over $k$ and all of whose prime ideals
are maximal ideals with residue field $k$. To see that the equality
holds, note that the number of morphisms $\overline{u}$ lying over
a given $u$ is equal to $[\kappa(u) : \kappa(v)]_s$ by
Fields, Lemma \href{https://stacks.math.columbia.edu/tag/09HJ}{lemma separable degree (Tag 09HJ)}.



\begin{definition}
\label{definition-weighting}
Let $f : X \to Y$ be a locally quasi-finite morphism. A
{\it weighting} or a {\it pond\'eration} of $f$ is a map
$w : X \to \mathbf{Z}$ such that for any diagram
$$
\xymatrix{
X \ar[d]_f & U \ar[l]^h \ar[d]^\pi \\
Y & V \ar[l]_g
}
$$
where $V \to Y$ is \'etale, $U \subset X_V$ is open, and $U \to V$ finite,
the function $\int_\pi (w \circ h)$ is locally constant.
\end{definition}

\begin{definition}
\label{fields-definition-insep-degree}
Let $E/F$ be an algebraic field extension. Let $E_{sep}$ be the subextension
found in Lemma \href{https://stacks.math.columbia.edu/tag/030K}{lemma separable first (Tag 030K)}.
\begin{enumerate}
\item The integer $[E_{sep} : F]$ is called the {\it separable
degree} of the extension. Notation $[E : F]_s$.
\item The integer $[E : E_{sep}]$ is called the {\it inseparable
degree}, or the {\it degree of inseparability} of the extension.
Notation $[E : F]_i$.
\end{enumerate}
\end{definition}

\begin{lemma}
\label{lemma-affineness-of-large-open}
Let $f : X \to Y$ be a morphism of affine schemes which is
quasi-finite and of finite presentation.
Let $w : X \to \mathbf{Z}_{> 0}$ be a positive weighting of $f$.
Let $d < \infty$ be the maximum value of $\int_f w$. The open
$$
Y_d = \{y \in Y \mid (\textstyle{\int}_f w)(y) = d \}
$$
of $Y$ is affine.
\end{lemma}

\begin{proof}
Observe that $\int_f w$ attains its maximum by Lemma \href{https://stacks.math.columbia.edu/tag/0F3K}{lemma max int w (Tag 0F3K)}.
The set $Y_d$ is open by Lemma \href{https://stacks.math.columbia.edu/tag/0F3J}{lemma semicontinuous int w (Tag 0F3J)}.
Thus the statement of the lemma makes sense.

\medskip\noindent
Reduction to the Noetherian case; please skip this paragraph.
Recall that a finite type morphism is quasi-finite if and only
if it has relative dimension $0$, see
Morphisms, Lemma \href{https://stacks.math.columbia.edu/tag/0397}{lemma locally quasi finite rel dimension 0 (Tag 0397)}.
By Lemma \href{https://stacks.math.columbia.edu/tag/05FJ}{lemma Noetherian approximation dimension d (Tag 05FJ)}
applied with $d = 0$ we can find a quasi-finite morphism $f_0 : X_0 \to Y_0$
of affine Noetherian schemes and a morphism $Y \to Y_0$ such that $f$
is the base change of $f_0$. Then we can write $Y = \lim Y_i$ as a directed
limit of affine schemes of finite type over $Y_0$, see
Algebra, Lemma \href{https://stacks.math.columbia.edu/tag/00QN}{lemma ring colimit fp (Tag 00QN)}.
By Lemma \href{https://stacks.math.columbia.edu/tag/0F3P}{lemma descend weighting (Tag 0F3P)}
we can find an $i$ such that our weighting $w$
descends to a weighting $w_i$ of the base change $f_i : X_i \to Y_i$
of $f_0$. Now if the lemma holds for $f_i, w_i$, then it implies
the lemma for $f$ as formation of $\int_f w$ commutes with base
change, see Lemma \href{https://stacks.math.columbia.edu/tag/0F39}{lemma weighting check after etale base change (Tag 0F39)}.

\medskip\noindent
Assume $X$ and $Y$ Noetherian. Let $X' \to Y'$ be the base change of $f$
by a morphism $g : Y' \to Y$. The formation of $\int_f w$ and hence
the open $Y_d$ commute with base change. If $g$ is finite and surjective, then
$Y'_d \to Y_d$ is finite and surjective. In this case proving that
$Y_d$ is affine is equivalent to showing that $Y'_d$ is affine, see
Cohomology of Schemes, Lemma
\href{https://stacks.math.columbia.edu/tag/01YQ}{lemma image affine finite morphism affine Noetherian (Tag 01YQ)}.

\medskip\noindent
We may choose an immersion $X \to T$ with $T$ finite over $Y$, see Lemma
\href{https://stacks.math.columbia.edu/tag/05K0}{lemma quasi finite separated pass through finite (Tag 05K0)}.
We are going to apply Morphisms, Lemma \href{https://stacks.math.columbia.edu/tag/03HW}{lemma massage finite (Tag 03HW)}
to the finite morphism $T \to Y$. This lemma tells us that
there is a finite surjective morphism $Y' \to Y$ such that
$Y' \times_Y T$ is a closed subscheme of a scheme $T'$ finite over $Y'$
which has a special form.
By the discussion in the first paragraph, we may replace $Y$ by $Y'$,
$T$ by $T'$, and $X$ by $Y' \times_Y X$.
Thus we may assume there is an immersion $X \to T$ (not necessarily
open or closed) and closed subschemes
$T_i \subset T$, $i = 1, \ldots, n$ where
\begin{enumerate}
\item $T \to Y$ is finite (and locally free),
\item $T_i \to Y$ is an isomorphism, and
\item $T = \bigcup_{i = 1, \ldots, n} T_i$ set theoretically.
\end{enumerate}
Let $Y' = \coprod Y_k$ be the disjoint union of the irreducible
components of $Y$ (viewed as integral closed subschemes of $Y$).
Then we may base change once more by $Y' \to Y$; here we
are using that $Y$ is Noetherian. Thus we may in
addition assume $Y$ is integral and Noetherian.

\medskip\noindent
We also may and do assume that $T_i \not = T_j$ if $i \not = j$ by
removing repeats. Since $Y$ and hence all $T_i$ are integral, this
means that if $T_i$ and $T_j$ intersect, then they intersect in a
closed subset which maps to a proper closed subset of $Y$.

\medskip\noindent
Observe that $V_i = X \cap T_i$ is a locally closed subset
which is in addition a closed subscheme of $X$ hence affine.
Let $\eta \in Y$ and $\eta_i \in T_i$ be the generic points.
If $\eta \not \in Y_d$, then $Y_d = \emptyset$ and we're done.
Assume $\eta \in Y_d$. Denote $I \in \{1, \ldots, n\}$
the subset of indices $i$ such that $\eta_i \in V_i$.
For $i \in I$ the locally closed subset $V_i \subset T_i$
contains the generic point of the irreducible space $T_i$
and hence is open. On the other hand, since $f$ is open
(Lemma \href{https://stacks.math.columbia.edu/tag/0F3C}{lemma weighting universally open (Tag 0F3C)}),
for any $x \in X$ we can find an $i \in I$ and a specialization
$\eta_i \leadsto x$. It follows that $x \in T_i$ and hence
$x \in V_i$. In other words, we see that $X = \bigcup_{i \in I} V_i$
set theoretically.
We claim that $Y_d = \bigcap_{i \in I} \Im(V_i \to Y)$; this will
finish the proof as the intersection of affine opens
$\Im(V_i \to Y)$ of $Y$ is affine.

\medskip\noindent
For $y \in Y$ let $f^{-1}(\{y\}) = \{x_1, \ldots, x_r\}$ in $X$.
For each $i \in I$ there is at most one $j(i) \in \{1, \ldots, x_r\}$
such that $\eta_i \leadsto x_{j(i)}$. In fact, $j(i)$ exists and is
equal to $j$ if and only if $x_j \in V_i$. If $i \in I$ is such that
$j = j(i)$ exists, then $V_i \to Y$ is an isomorphism in a neighbourhood
of $x_j \mapsto y$. Hence $\bigcup_{i \in I,\ j(i) = j} V_i \to Y$
is finite after replacing source and target by neighbourhoods of
$x_j \mapsto y$. Thus the definition of a weighting tells us that
$w(x_j) = \sum_{i \in I,\ j(i) = j} w(\eta_i)$.
Thus we see that
$$
(\textstyle{\int}_f w)(\eta) =
\sum\nolimits_{i \in I} w(\eta_i) \geq
\sum\nolimits_{j(i)\text{ exists}} w(\eta_i) =
\sum\nolimits_j w(x_j) = (\textstyle{\int}_f w)(y)
$$
Thus equality holds if and only if $y$ is contained in
$\bigcap_{i \in I} \Im(V_i \to Y)$ which is what we wanted to show.
\end{proof}
\end{document}
